\section{The Architectures Under Study}
\label{sec:architectures}

We compare thirteen architectures: a no-retrieval floor, the naive
baseline at three context budgets, and nine refinements drawn from the
families of Section~\ref{sec:related}. Figure~\ref{fig:taxonomy} groups
them by the shape of their control flow, following the taxonomy of
\cite{jin2024flashrag}: \emph{sequential} architectures run the pipeline of
Figure~\ref{fig:naive-pipeline} once, changing only which passages reach
the generator; \emph{branching} architectures add a generation step before
retrieval; \emph{conditional} architectures inspect intermediate state and
choose among actions; \emph{iterative} architectures loop.
Table~\ref{tab:qualitative} summarises what each one changes and what it
costs. Each architecture below is described with the same template ---
what it changes, how it works step by step, and what it costs --- and each
description states the measured failure mode it targets, anticipating the
diagnosis of Section~\ref{sec:results-diagnosis}.

\begin{figure}[t]
\centering
\begin{tikzpicture}[node distance=3mm and 5mm]
\node[comp, fill=black!10, minimum width=32mm] (root) {13 architectures};
\node[comp, fill=blue!10, below left=8mm and 34mm of root, minimum width=30mm] (seq) {Sequential\\ \scriptsize one pass, better passages};
\node[comp, fill=violet!10, below left=8mm and -6mm of root, minimum width=26mm] (bra) {Branching\\ \scriptsize generate, then search};
\node[comp, fill=yellow!20, below right=8mm and -6mm of root, minimum width=26mm] (con) {Conditional\\ \scriptsize judge, then act};
\node[comp, fill=orange!15, below right=8mm and 30mm of root, minimum width=24mm] (ite) {Iterative\\ \scriptsize loop};
\draw[arr] (root) -- (seq); \draw[arr] (root) -- (bra);
\draw[arr] (root) -- (con); \draw[arr] (root) -- (ite);
\node[lbl, below=2mm of seq, text width=34mm] {ClosedBook$^{*}$\\ NaiveRAG (5/10/12)\\ BM25 \quad Hybrid\\ Rerank \quad Neighbour\\ \textbf{Fused} (\S\ref{sec:cascade})};
\node[lbl, below=2mm of bra, text width=26mm] {HyDE};
\node[lbl, below=2mm of con, text width=30mm] {CRAG (5/10)\\ Adaptive-RAG\\ \textbf{Cascade} (\S\ref{sec:cascade})};
\node[lbl, below=2mm of ite, text width=24mm] {IRCoT\\ (zero- and one-shot)};
\end{tikzpicture}
\caption{The compared architectures grouped by control flow.
$^{*}$ClosedBook retrieves nothing and serves as the floor. The two
architectures in bold are introduced in Section~\ref{sec:cascade}.}
\label{fig:taxonomy}
\end{figure}

\begin{table}[t]
\centering\small
\caption{Qualitative comparison. ``Calls'' is generator calls per question;
``budget'' is passages placed in the generator's context. The measured
per-question costs are in Table~\ref{tab:cost}.}
\label{tab:qualitative}
\begin{tabular}{lllll}
\toprule
Architecture & Modifies & Calls & Budget & Targets \\
\midrule
ClosedBook & removes retrieval & 1 & 0 & (cost floor) \\
NaiveRAG & --- (baseline) & 1 & 5 & --- \\
NaiveRAG 10/12 & context budget & 1 & 10/12 & (budget controls) \\
BM25 & retriever (lexical) & 1 & 5 & rare surface forms \\
Hybrid & retriever (dense+lexical) & 1 & 5 & rare surface forms \\
Rerank & ranking (cross-encoder) & 1 & 5 & right article ranked low \\
Neighbour & chunking (expand hits) & 1 & $\sim$11.5 & wrong slice of article \\
HyDE & query (drafted passage) & 2 & 5 & unreachable by query \\
CRAG & adds retrieval judge & 3 & 5/10 & undetected bad retrieval \\
IRCoT & query (iterative reasoning) & 2--6 & $\le$15 & multi-hop questions \\
Adaptive-RAG & routes between strategies & 2--7 & 0--15 & heterogeneous difficulty \\
\midrule
Cascade (proposed) & retrieval + escalation & $\sim$1.3 & $\le$5 & all of the above \\
\bottomrule
\end{tabular}
\end{table}

\subsection{ClosedBook: the floor}
\label{sec:arch-closedbook}

The generator answers from parametric knowledge alone, with the same
answer-only prompt as every other row minus the documents. This is the
floor each retrieval number must be read against: a 7B model has memorised
a large amount of Wikipedia-scale trivia (it answers
\ClosedBookTriviaQAEM{} of TriviaQA without any retrieval), so a high RAG
score does not by itself show that retrieval contributed anything. The
difference between ClosedBook and NaiveRAG is what retrieval is actually
worth on each question set --- and on one set that difference will turn
out to be an artefact of abstention behaviour rather than knowledge
(Section~\ref{sec:results-worth}).

\subsection{NaiveRAG: the baseline and its budget controls}
\label{sec:arch-naive}

The pipeline of Figure~\ref{fig:naive-pipeline} verbatim: embed the
question, retrieve the top five passages, generate. Two further rows run
the identical pipeline with ten and twelve passages. These exist because
several refinements below place more than five passages in the context, and
context size alone is worth measurable Exact Match (naive-10 gains $+0.022$
over naive-5 on 2WikiMultihopQA); comparing an architecture that reads ten
passages against a five-passage baseline would credit the architecture with
what is merely a larger budget. Throughout the article, every comparison
is made at matched budget.

\subsection{BM25 and Hybrid retrieval}
\label{sec:arch-hybrid}

\begin{figure}[t]
\centering
\begin{tikzpicture}[node distance=4mm and 7mm]
\node[q] (q) {Question};
\node[ret, above right=2mm and 9mm of q] (dense) {Dense search\\ \scriptsize bi-encoder, top 50};
\node[ret, below right=2mm and 9mm of q] (lex) {Lexical search\\ \scriptsize BM25, top 50};
\node[comp, fill=blue!6, right=30mm of q] (rrf) {RRF\\ \scriptsize merge ranks};
\node[llm, right=of rrf] (gen) {Generator};
\node[ans, right=of gen] (a) {Answer};
\draw[arr] (q) -- (dense); \draw[arr] (q) -- (lex);
\draw[arr] (dense) -- (rrf); \draw[arr] (lex) -- (rrf);
\draw[arr] (rrf) -- node[lbl, above] {top 5} (gen);
\draw[arr] (gen) -- (a);
\end{tikzpicture}
\caption{Hybrid retrieval. Dense and lexical rankings over the same corpus
are merged by Reciprocal Rank Fusion; the BM25 row is this diagram with
the dense branch removed.}
\label{fig:arch-hybrid}
\end{figure}

The two retrievers fail in opposite ways. The dense retriever understands
paraphrase but blurs rare surface forms (the lossy compression of
Section~\ref{sec:embeddings}); the lexical one matches names exactly but is
blind to a question that never uses the passage's vocabulary. Hybrid
retrieval runs both and merges the rankings with Reciprocal Rank
Fusion~\cite{cormack2009rrf}:
\begin{equation}
\mathrm{RRF}(d) \;=\; \sum_{r \,\in\, \text{rankings}} \frac{1}{k + r(d)},
\qquad k = 60,
\label{eq:rrf}
\end{equation}
where $r(d)$ is the rank of document $d$ in each list. RRF combines
\emph{ranks}, not scores, which is what makes the combination possible at
all: a cosine similarity and a BM25 score live on incomparable scales,
while their ranks do not. The BM25-only row is the same code with the dense
branch disabled, so the difference between the two rows measures exactly
``was the dense list fused in'' and nothing else. Cost: one extra lexical
search on the CPU; the generator is still called once.

\subsection{Rerank: cross-encoder reordering}
\label{sec:arch-rerank}

\begin{figure}[t]
\centering
\begin{tikzpicture}[node distance=4mm and 7mm]
\node[q] (q) {Question};
\node[ret, right=of q] (dense) {Dense search\\ \scriptsize top 100};
\node[ret, right=of dense, fill=red!12] (ce) {Cross-encoder\\ \scriptsize scores all 100};
\node[llm, right=of ce] (gen) {Generator};
\node[ans, right=of gen] (a) {Answer};
\draw[arr] (q) -- (dense);
\draw[arr] (dense) -- (ce);
\draw[arr] (ce) -- node[lbl, above] {top 5} (gen);
\draw[arr] (gen) -- (a);
\end{tikzpicture}
\caption{Cross-encoder reranking: a deep bi-encoder shortlist is reordered
by a jointly-encoding relevance model; only the reordered top five reach
the generator.}
\label{fig:arch-rerank}
\end{figure}

A deep shortlist (100 candidates) is retrieved with the fast bi-encoder and
reordered with a cross-encoder (Section~\ref{sec:encoders}); the top five
survivors reach the generator. The target failure is ``right article
ranked too low'': for a third of the diagnosed failures on 2WikiMultihopQA
the answer-bearing article is already in the ranking, just below the
cut-off, and nothing that reasons over the top five can recover it --- the
fix has to change which five are chosen. Unlike the iterative methods this
adds no generator call; the cost is one cross-encoder pass over the
shortlist.

\subsection{Neighbour expansion}
\label{sec:arch-neighbour}

Wikipedia articles are chunked into $\sim$100-word passages with no regard
for where a fact ends, so the retriever can land on the paragraph about the
right entity while the sentence carrying the answer sits in the next chunk.
Neighbour expansion retrieves as NaiveRAG does, then adds, for every hit,
the passages adjacent to it in the same article (radius one, capped at
fifteen passages, $\sim$11.5 in practice). It is included to \emph{test}
the chunking hypothesis rather than because our diagnosis supports it ---
the ``wrong slice'' bucket is only 5--14\% of failures --- and it is read
against the twelve-passage control, because its expanded context would
otherwise masquerade as a retrieval gain.

\subsection{HyDE: search with a hypothetical document}
\label{sec:arch-hyde}

\begin{figure}[t]
\centering
\begin{tikzpicture}[node distance=4mm and 7mm]
\node[q] (q) {Question};
\node[llm, right=of q] (draft) {Generator\\ \scriptsize drafts a passage};
\node[ret, right=of draft] (enc) {Embed draft\\ $+$ question};
\node[ret, right=of enc] (search) {Dense search\\ \scriptsize mean vector};
\node[llm, right=of search] (gen) {Generator};
\node[ans, below=of gen] (a) {Answer};
\draw[arr] (q) -- (draft);
\draw[arr] (draft) -- (enc);
\draw[arr] (enc) -- (search);
\draw[arr] (search) -- node[lbl, above] {top 5} (gen);
\draw[arr] (gen) -- (a);
\end{tikzpicture}
\caption{HyDE. The generator first writes the encyclopedia passage it
imagines would answer the question; the search runs on the average of the
draft's and the question's embeddings. The draft's factual errors are
discarded --- only its vector is used.}
\label{fig:arch-hyde}
\end{figure}

A question and the passage that answers it are different kinds of text:
``Where was the place of death of the director of film Beat Girl?'' shares
almost no vocabulary with a paragraph about Edmond T.\ Gr\'eville. HyDE
\cite{gao2022hyde} removes the mismatch by asking the generator to write
the passage it \emph{thinks} would answer the question and searching with
that instead. The draft is usually wrong in its details and it does not
matter: it only has to look like the right passage so that the encoder
places it near the real one. The question's own embedding is averaged in
as a safety net, so an off-topic draft degrades towards NaiveRAG instead
of chasing a hallucination. This targets the ``unreachable by this query''
failure bucket --- 53\% of failures on MuSiQue. Cost: one extra generator
call, the cheapest of the query-side methods.

\subsection{IRCoT: interleaved retrieval and reasoning}
\label{sec:arch-ircot}

\begin{figure}[t]
\centering
\begin{tikzpicture}[node distance=4mm and 6mm]
\node[q] (q) {Question};
\node[ret, right=of q] (r0) {Retrieve\\ \scriptsize top 5};
\node[llm, right=of r0] (step) {Reason\\ \scriptsize one sentence};
\node[dec, right=of step] (done) {answer\\found?};
\node[llm, right=of done] (gen) {Generator};
\node[ans, right=of gen] (a) {Answer};
\draw[arr] (q) -- (r0);
\draw[arr] (r0) -- (step);
\draw[arr] (step) -- (done);
\draw[arr] (done) -- node[lbl, above] {yes} (gen);
\draw[arr] (done.south) to[out=-90,in=-90] node[lbl, below] {no: sentence becomes next query} (r0.south);
\draw[arr] (gen) -- (a);
\end{tikzpicture}
\caption{IRCoT. Each reasoning sentence names entities the question never
mentioned, and is used verbatim as the next retrieval query; accumulated
passages (up to 15) go to the same reader prompt as NaiveRAG.}
\label{fig:arch-ircot}
\end{figure}

A single query built from the question alone is structurally insufficient
for multi-hop questions: ``What is the capital of the country where X was
born?'' never names the country, so no single query can reach the passage
that answers it. Measured on the baseline, recall over the context falls
from \NaiveRAGHotpotQARecallInContext{} on HotpotQA to
\NaiveRAGTwoWikiRecallInContext{} on 2WikiMultihopQA to
\NaiveRAGMuSiQueRecallInContext{} on MuSiQue as the number of required hops
grows. IRCoT~\cite{trivedi2023ircot} alternates: reason one sentence over
what has been retrieved, use that sentence as the next query, add what
comes back, repeat (up to four rounds or fifteen passages). The final
answer is produced by the same reader prompt as NaiveRAG over the
accumulated passages, so the only difference between the two systems is
the set of passages in context. We run a zero-shot variant and, because
the original method prompts the reasoning step with a worked example, a
one-shot variant with a single hand-written two-hop demonstration ---
identical for every question set --- as a separate row, so that a
under-prompted-baseline objection can be answered with data. Cost: two to
six generator calls per question.

\subsection{CRAG: corrective retrieval (offline, corpus-only)}
\label{sec:arch-crag}

\begin{figure}[t]
\centering
\begin{tikzpicture}[node distance=4mm and 6mm]
\node[q] (q) {Question};
\node[ret, right=of q] (r0) {Retrieve\\ \scriptsize top 5};
\node[llm, right=of r0] (grade) {Grade\\ \scriptsize which passages help?};
\node[dec, right=of grade] (branch) {verdict};
\node[llm, above right=1mm and 8mm of branch] (rewrite) {Rewrite query,\\re-search};
\node[llm, below right=1mm and 8mm of branch] (gen) {Generator};
\node[ans, right=of gen] (a) {Answer};
\draw[arr] (q) -- (r0);
\draw[arr] (r0) -- (grade);
\draw[arr] (grade) -- (branch);
\draw[arr] (branch) -- node[lbl, above left] {none/some relevant} (rewrite);
\draw[arr] (branch) -- node[lbl, below left] {all relevant} (gen);
\draw[arr] (rewrite.east) to[out=0,in=60] node[lbl, right] {replaced/merged passages} (gen.north east);
\draw[arr] (gen) -- (a);
\end{tikzpicture}
\caption{CRAG (offline form). The generator grades its own retrieval; on a
poor grade the query is rewritten and the corpus searched again, replacing
(no passage relevant) or topping up (some relevant) the context.}
\label{fig:arch-crag}
\end{figure}

NaiveRAG hands the generator whatever came back and hopes it is relevant;
the baseline measurements say how often that hope is misplaced (a gold
answer appears in the retrieved context for only
\NaiveRAGTwoWikiAnswerInContext{} of 2WikiMultihopQA questions).
CRAG~\cite{yan2024crag} adds the missing step: judge the retrieval before
using it, and act when it is bad. Our implementation deviates from the
paper in three deliberate, named ways, and is labelled ``CRAG (offline,
corpus-only)'' throughout: the generator itself grades relevance in one
call (the paper trains a T5 evaluator; a trained evaluator would confound
the controlled comparison with an extra model); the corrective action is a
query rewrite over the same corpus (the paper uses web search, which would
change the available knowledge, not the architecture); and the paper's
knowledge-refinement step is absent. A ten-passage variant pairs with the
matched-budget control. Cost: three generator calls per question.

\subsection{Adaptive-RAG: routing by predicted difficulty}
\label{sec:arch-adaptive}

\begin{figure}[t]
\centering
\begin{tikzpicture}[node distance=4mm and 7mm]
\node[q] (q) {Question};
\node[llm, right=of q] (route) {Classify\\ \scriptsize A / B / C};
\node[comp, fill=black!5, above right=6mm and 16mm of route] (cb) {ClosedBook};
\node[comp, fill=black!5, right=16mm of route] (naive) {NaiveRAG};
\node[comp, fill=black!5, below right=6mm and 16mm of route] (ircot) {IRCoT};
\node[ans, right=42mm of route] (a) {Answer};
\draw[arr] (q) -- (route);
\draw[arr] (route) -- node[lbl, above, sloped] {famous fact} (cb);
\draw[arr] (route) -- node[lbl, above, pos=0.45] {one lookup} (naive);
\draw[arr] (route) -- node[lbl, below, sloped] {multi-hop} (ircot);
\draw[arr] (cb) -- (a); \draw[arr] (naive) -- (a); \draw[arr] (ircot) -- (a);
\end{tikzpicture}
\caption{Adaptive-RAG (training-free form). One classification call routes
each question to the no-retrieval, single-step, or iterative branch; the
branches are the very architectures of the neighbouring rows.}
\label{fig:arch-adaptive}
\end{figure}

Every other architecture treats all questions alike, and the measurements
say that is wrong in both directions: on TriviaQA the generator already
knows \ClosedBookTriviaQAEM{} of the answers with no retrieval, so running
IRCoT over them spends up to six calls to arrive where one would have
done, while on MuSiQue a single query reaches the answer-bearing article
for barely a fifth of questions. Adaptive-RAG~\cite{jeong2024adaptiverag}
routes each question --- here with a single generator classification call
rather than the paper's trained classifier, for the same
no-extra-models reason as CRAG --- to closed-book, single-step, or IRCoT.
The router's branches are built from the same components as the
architectures they are compared against, so its row is comparable with
theirs by construction. Cost: one call plus whatever the chosen branch
costs (2.1--3.0 in practice).
